\documentclass[letterpaper,11pt]{article}

\usepackage[letterpaper,left=0.80in,right=0.80in,top=0.70in,bottom=0.70in]{geometry}
\usepackage[hidelinks]{hyperref}
\usepackage[english]{babel}
\input{glyphtounicode}

\pagestyle{empty}
\urlstyle{same}
\raggedright
\setlength{\parindent}{0pt}
\setlength{\parskip}{8pt}
\pdfgentounicode=1

\begin{document}

\begin{center}
  {\LARGE \textbf{Aryan Miriyala}}\\[-1pt]
  \small
  +1 419-315-0444 $|$ \href{mailto:aryanmiriyala@gmail.com}{aryanmiriyala@gmail.com} $|$ \href{https://www.linkedin.com/in/aryan-miriyala-7788a21b6/}{linkedin.com/in/aryan-miriyala}
\end{center}

October 15, 2026

Deloitte Hiring Team\\
Deloitte Consulting LLP

Dear Deloitte Hiring Team,

I am applying for the Agentic Software Engineer I position on the Customer team. What draws me to the role is the specific problem it names: turning ambiguous requirements into structured context that an AI coding agent can act on, then reviewing what comes back against the original requirements. That is the part of agent-assisted delivery that determines whether a workflow produces something a team can ship, and it is the work I have spent my recent internships doing.

Most of my day at SmartSolve Industries was spent inside that loop. I drove development with Claude Code, GitHub Copilot, Cursor, and Codex across planning, debugging, code review, and self-review, and my responsibility was to validate agent-generated code for quality and alignment to requirements before it merged. The harder part was upstream of that. Operations requests arrived as one-sentence problems, so I structured them into technical tasks and explicit context for the agent to work from, and wrote the plan so a reviewer could check the output against what had actually been asked. That work turned into a quality-management dependency register that both technical and non-technical reviewers could trace.

At Actual Reality Technologies I built the release controls on the other side of the same loop. I deployed a LiteLLM AI Gateway that routes model calls by policy across several AI projects, then added guardrails, Langfuse tracing, Slack alerts, and 44 red-team fixtures, along with CI and deploy gates that block unsafe routing configurations before they ship. I also built a remote Model Context Protocol server with keyed access so customer and admin AI tools can query ticket and project status under scoped security and cost controls. What I take from that is that an agentic workflow only earns trust when the delivery controls around it are as carefully built as the model calls inside it.

Earlier, as a temporary Product Owner for an internal application at the American Association of Insurance Services, I translated business requirements into implementation tasks and coordinated delivery priorities, which is the same requirements-decomposition skill the role calls for. I hold an M.S. in Computer Science from Bowling Green State University with a 4.00 GPA, alongside earlier internship work across AWS data platforms, access control, and automated testing.

I would welcome the chance to discuss how this experience contributes to the Customer team's AI-assisted delivery work.

Sincerely,

Aryan Miriyala

\end{document}
